\setcounter{page}{1}
\renewcommand{\thepage}{S\arabic{page}}
\setcounter{section}{0}
\setcounter{equation}{0}
\setcounter{table}{0}
\setcounter{figure}{0}
\setcounter{proposition}{0}
\setcounter{corollary}{0}
\setcounter{assumption}{0}
\setcounter{definition}{0}
\renewcommand{\thesection}{S\arabic{section}}
\renewcommand{\theequation}{S\arabic{equation}}
\renewcommand{\thetable}{S\arabic{table}}
\renewcommand{\thefigure}{S\arabic{figure}}
\renewcommand{\theproposition}{S\arabic{proposition}}
\renewcommand{\thecorollary}{S\arabic{corollary}}
\renewcommand{\theassumption}{S\arabic{assumption}}
\renewcommand{\thedefinition}{S\arabic{definition}}
\renewcommand{\theHsection}{S\arabic{section}}
\renewcommand{\theHequation}{S\arabic{equation}}
\renewcommand{\theHtable}{S\arabic{table}}
\renewcommand{\theHfigure}{S\arabic{figure}}
\renewcommand{\theHproposition}{S\arabic{proposition}}
\providecommand{\theHcorollary}{\arabic{corollary}}
\providecommand{\theHassumption}{\arabic{assumption}}
\providecommand{\theHdefinition}{\arabic{definition}}
\renewcommand{\theHcorollary}{S\arabic{corollary}}
\renewcommand{\theHassumption}{S\arabic{assumption}}
\renewcommand{\theHdefinition}{S\arabic{definition}}
\allowdisplaybreaks[0]
\makeatletter
\let\SIoriginaladdtocontents\addtocontents
\renewcommand{\addtocontents}[2]{
  \def\SIextension{#1}
  \def\SItoc{toc}
  \ifx\SIextension\SItoc
    \SIoriginaladdtocontents{sitoc}{#2}
  \else
    \SIoriginaladdtocontents{#1}{#2}
  \fi
}
\renewcommand*{\l@section}[2]{\begingroup\bfseries\@dottedtocline{1}{0em}{2.2em}{#1}{#2}\endgroup}

\begingroup
\def\@title{Supplementary Information for\\``Manuscript Title''}
\def\@author{Author One \and Author Two}
\def\@date{}
\let\and\SIoriginaland
\SIoriginalmaketitle
\endgroup
\thispagestyle{plain}
\makeatother
{\setcounter{tocdepth}{1}\makeatletter
\let\SIoriginalstarttoc\@starttoc
\renewcommand{\@starttoc}[1]{\SIoriginalstarttoc{sitoc}}
\tableofcontents
\makeatother}
\clearpage
\onehalfspacing

\section{Proof of the main result}\label{SI-sec:proof}

This section proves Proposition~\ref{prop:mean} from Section~\ref{sec:framework} of the main text.

\begin{proof}[Proof of Proposition~\ref{prop:mean}]
By Assumption~\ref{ass:sampling} and Equation~\eqref{eq:mean},
\begin{align}
\E[\bar Y]
&=\frac{1}{n}\sum_{i=1}^{n}\E[Y_i]
=\mu,\label{SI-eq:expectation}\\
\Var(\bar Y)
&=\frac{1}{n^2}\sum_{i=1}^{n}\Var(Y_i)
=\frac{\sigma^2}{n}.\label{SI-eq:variance}
\end{align}
Independence makes the covariance terms zero. These equalities establish Equation~\eqref{eq:moments}.
\end{proof}

\section{A weighted extension}\label{SI-sec:extension}

\subsection{Conditions and result}\label{SI-sec:weighted-result}

\begin{assumption}[Fixed weights]\label{SI-ass:weights}
The nonrandom weights $w_1,\ldots,w_n$ are nonnegative and satisfy $\sum_{i=1}^{n}w_i=1$.
\end{assumption}

\begin{definition}[Weighted mean]\label{SI-def:weighted}
The weighted mean is $\sum_{i=1}^{n}w_iY_i$.
\end{definition}

\begin{proposition}[Weighted moments]\label{SI-prop:weighted}
Under Assumptions~\ref{ass:sampling} and~\ref{SI-ass:weights},
\begin{equation}\label{SI-eq:weighted-moments}
\E\left[\sum_{i=1}^{n}w_iY_i\right]=\mu,
\qquad
\Var\left(\sum_{i=1}^{n}w_iY_i\right)=\sigma^2\sum_{i=1}^{n}w_i^2.
\end{equation}
\end{proposition}

\begin{proof}
Linearity of expectation and $\sum_i w_i=1$ give the mean. Independence gives the variance by the same argument as in Equation~\eqref{SI-eq:variance}.
\end{proof}

\begin{corollary}[Equal weights]\label{SI-cor:equal}
Setting $w_i=1/n$ in Proposition~\ref{SI-prop:weighted} recovers Proposition~\ref{prop:mean}.
\end{corollary}

\clearpage
\subsection{Illustration}\label{SI-sec:illustration}

Table~\ref{SI-tab:weights} gives two weight choices. Figure~\ref{SI-fig:weights} illustrates the construction for two observations, using the same shared drawing styles as Figure~\ref{fig:diagram} in the main text.

\begin{table}[!htbp]
\centering
\caption{Illustrative fixed weights.}
\label{SI-tab:weights}
\begin{tabularx}{\linewidth}{@{}Xll@{}}
\toprule
Choice & Weights & Variance \\
\midrule
Equal weights & $w_i=1/n$ & $\sigma^2/n$ \\
First observation only & $w_1=1$, $w_i=0$ for $i>1$ & $\sigma^2$ \\
\bottomrule
\end{tabularx}
\end{table}

\begin{figure}[!htbp]
\centering
\begin{tikzpicture}[dag]
\node[dag observed,label=left:{$Y_1$}] (Y1) at (0,1) {};
\node[dag observed,label=left:{$Y_2$}] (Y2) at (0,-1) {};
\node[dag observed,label=right:{$w_1Y_1+w_2Y_2$}] (mean) at (3,0) {};
\draw[dag arrow] (Y1) -- (mean);
\draw[dag arrow] (Y2) -- (mean);
\end{tikzpicture}

\caption{Two observations contribute to a weighted mean. Arrows indicate the construction, not causal effects.}
\label{SI-fig:weights}
\end{figure}

\FloatBarrier
